\FloatBarrier
\newpage
\section{概率模型}
\label{sec:prob-models}

计数确定可能结果的个数，概率还需要说明各结果怎样分配权重。先把一次试验记录清楚，才能判断所用的计数是否适合当前概率计算。

\subsection{样本空间}

\begin{definition}
一次随机试验的一个完整可能结果称为{\heiti 样本点}，全体样本点组成的集合称为{\heiti 样本空间}，通常记为\(\Omega\)。
\end{definition}

例如，连续抛掷两次硬币，分别用\(H,T\)记录正反面，可以取
\[
 \Omega=\{HH,HT,TH,TT\}.
\]
若只关心正面次数，也可以记录为\(\{0,1,2\}\)，但两个模型保留的信息不同。前者可以区分第一次正面还是第二次正面，后者把\(HT,TH\)合成了同一种数值记录。三个数值并不会因此自动成为等可能结果。

本章先研究有限样本空间。建立它时，应使每次试验恰能归入一个样本点，所有可能结果均在其中。

\subsection{事件}

\begin{definition}
样本空间的一个子集称为{\heiti 事件}。试验结果属于这个子集时，就说该事件发生。\(\Omega\)称为必然事件，\(\varnothing\)称为不可能事件。
\end{definition}

事件\(A\cup B\)表示\(A,B\)至少一个发生，\(A\cap B\)表示两者同时发生。补事件\(A^c=\Omega\setminus A\)表示\(A\)不发生；\(A\setminus B=A\cap B^c\)表示\(A\)发生而\(B\)不发生。

若\(A\cap B=\varnothing\)，称\(A,B\)互斥。若\(A\subseteq B\)，则\(A\)发生必有\(B\)发生。例如在两次抛币中，“两次正面”包含于“至少一次正面”。事件的关系由它所含的样本点确定，可以直接使用集合运算。

例如两次抛币，事件\(A=\{HH,HT\}\)表示第一次正面，\(B=\{HH,TH\}\)表示第二次正面。“恰一次正面”便是\((A\setminus B)\cup(B\setminus A)=\{HT,TH\}\)。

若\(\Omega\)有\(N\)个样本点，事件共有\(2^N\)个：每个样本点都可以选择属于或不属于所取的子集。因此，样本点的个数与事件的个数是不同的计数对象。

\begin{example}
从编号1至5的球中依次不放回取两球。分别用有序对和无序集合表示样本空间，并判断哪种记录能表示事件“第一次取到1号球”。
\end{example}

\newpage
\subsection{概率分配}

对有限样本空间\(\Omega=\{\omega_1,\ldots,\omega_N\}\)，给各样本点分配数\(p_i\)，满足
\[
 p_i\ge0,\qquad \sum_{i=1}^N p_i=1.
\]
定义事件\(A\)的{\heiti 概率}为
\begin{equation}\label{eq:prob-finite-probability}
 P(A)=\sum_{\omega_i\in A}p_i.
\end{equation}
于是\(P(\Omega)=1\)、\(P(\varnothing)=0\)。互斥事件的并，其概率等于概率之和。由集合分解还得到
\begin{align}
 P(A^c)&=1-P(A),\label{eq:prob-complement}\\
 P(A\cup B)&=P(A)+P(B)-P(A\cap B).\label{eq:prob-addition}
\end{align}
若\(A\subseteq B\)，由\(B=A\cup(B\setminus A)\)得到\(P(A)\le P(B)\)。这些性质都来自非负权重的求和。

例如两次独立公平抛币的四个序列各分配\(1/4\)，正面次数\(0,1,2\)的概率分别为\(1/4,1/2,1/4\)。记录合并时，需要把相应样本点的概率相加。

样本空间、其事件集合及概率分配共同组成一个{\heiti 有限概率空间}。机制不同，可以在同一结果集合上给出不同的概率分配。

\subsection{古典概型}

若有限样本空间中各样本点等可能，则每个样本点的概率为\(1/N\)，式\eqref{eq:prob-finite-probability}成为
\begin{equation}\label{eq:prob-classical}
 P(A)=\frac{|A|}{|\Omega|}.
\end{equation}
这种模型称为{\heiti 古典概型}。加法、乘法原理与排列组合公式在这里提供分子、分母的计数。

例如从\(n\)个不同对象中均匀选取\(r\)个，所有\(r\)元子集等可能；依次不放回抽取时，所有长度为\(r\)的无重复序列也等可能。每个子集都对应\(r!\)个序列，因此两种记录下，同一个事件的概率相同。计算时只需前后一致地选择一种记录。

\begin{example}[2023年全国甲卷文科，题面整理]
某校文艺部有4名学生，高一、高二各2名。从中随机选2名组织文艺汇演，求两人来自不同年级的概率。
\end{example}

\begin{example}[2024年全国甲卷文科，题面整理]
甲、乙、丙、丁各出场一次，次序由随机抽签确定。求丙不是第一个出场，且甲或乙最后出场的概率。
\end{example}

\newpage
\subsection{频率与模拟}

在\(n\)次试验中，事件\(A\)发生\(N_A\)次，称\(N_A/n\)为\(A\)在这组试验中的{\heiti 频率}。它由实际记录得到；模型中的\(P(A)\)则由概率分配确定。若模型含未知参数，可以用观测频率估计它。

随机模拟可以给出新的试验记录。例如用公平随机整数\(0,1,\ldots,9\)，将\(0,1,2,3\)记为成功，便得到成功概率为\(0.4\)的试验。独立重复这个步骤，记录不同试验次数下的成功频率，就能观察波动如何变化。

\begin{figure}[!ht]
\centering
\includegraphics[width=.78\linewidth]{probability-frequency}
\caption{成功概率为\(0.4\)的独立重复试验的一条模拟记录。虚线表示模型概率。}
\label{fig:prob-frequency}
\end{figure}

模拟中的频率仍会波动。为什么较大的试验次数通常能带来更可靠的估计，将由后面的方差与大数定律说明。

\begin{exercise}
\mbox{}\par\nobreak

（1）独立抛掷两枚公平骰子。分别用有序点数对与点数和记录结果，说明哪一种记录的样本点等可能，并求点数和为7的概率。

（2）设\(P(A)=0.6\)、\(P(B)=0.5\)、\(P(A\cup B)=0.7\)，求\(P(A\cap B)\)、\(P(A\setminus B)\)和\(P(A^c\cap B^c)\)。

（3）从编号1至6的球中不放回抽取3个。分别按有序抽取和无序子集计数，求其中恰有两个偶数的概率。

（4）某次模拟做1000次试验，事件发生397次。再做1000次时，是否必须恰好发生403次，才能使总频率等于模型给定的\(0.4\)？说明模型概率与实际记录各自表示什么。
\end{exercise}
